\documentclass[12pt,a4paper,oneside]{article}

\usepackage[spanish]{babel}

\def\mydate{\leavevmode\hbox{\the\year-\twodigits\month-\twodigits\day}}
\def\twodigits#1{\ifnum#1<10 0\fi\the#1}

\usepackage{fancyhdr}
\usepackage{geometry}
\usepackage{graphicx}
\usepackage{wrapfig}
\usepackage{lastpage}
\usepackage[hidelinks]{hyperref}
\usepackage{authblk}
\usepackage{bookmark}

\usepackage[utf8]{inputenc} % Required for inputting international characters
\usepackage[T1]{fontenc} % Output font encoding for international characters
\usepackage{mathpazo} % Palatino font

\makeatletter
\newcommand{\subtitledoc}[1]{\newcommand{\@subtitledoc}{#1}}
\newcommand{\instituto}[1]{\newcommand{\@instituto}{#1}}
\newcommand{\carrera}[1]{\newcommand{\@carrera}{#1}}
\newcommand{\professor}[1]{\newcommand{\@professor}{#1}}
\newcommand{\catedra}[1]{\newcommand{\@catedra}{#1}}
\newcommand{\curso}[1]{\newcommand{\@curso}{#1}}
\newcommand{\legajo}[1]{\newcommand{\@legajo}{#1}}
\newcommand{\footerauthor}[1]{\newcommand{\@footerauthor}{#1}}
\newcommand{\footerlegajo}[1]{\newcommand{\@footerlegajo}{#1}}
\newcommand{\footercatedra}[1]{\newcommand{\@footercatedra}{#1}}

%Configuracion de hoja (margenes y tamaño)
\geometry{a4paper,margin=1in}
\setlength\headheight{28pt}

%Metadata para pdf
%\hypersetup{
%    pdftitle={\@title},
%    pdfsubject={\@catedra\ - \@subtitledoc},
%    pdfauthor={\@author\ - \@legajo\ - \@curso},
%    pdfkeywords={\@title\ \@catedra\ \@subtitledoc\ \@author\ \@legajo\ \@curso}
%}

%formato de encabezado y pie para todas las paginas.
\fancyhead[L]{
  \begin{minipage}[b]{7.5mm}
    \includegraphics[width=7mm]{imagenes/COPYlogo-utn.eps}
  \end{minipage}
  \begin{minipage}[b]{85mm}
    \textbf{Alumnos: }\@footerauthor\\
    \textbf{Legajos: }\@footerlegajo
  \end{minipage}
}
\fancyhead[R]{
  %\textbf{Comisión:} \@curso 
  \textbf{Fecha:} \mydate \\
  \textbf{Cátedra:} \@footercatedra \text{ (}\@curso\text{)}
}
\fancyfoot[C]{} %eliminar antiguo numero de pagina
\fancyfoot[R]{Página \thepage\ de \pageref{LastPage}}
\renewcommand{\headrulewidth}{0.5pt}
\renewcommand{\footrulewidth}{0.5pt}
\pagestyle{fancy}

\addto\captionsspanish{%
  \renewcommand{\contentsname}%
  {CONTENIDO}%
}

\renewcommand{\maketitle}{%
  \newpage
  \thispagestyle{empty}

  \begin{center}

    \textsc{\LARGE \@instituto}\par
    \vspace{0.5cm}
    \textsc{\Large \@carrera}\par
    \vspace{1cm}
    \includegraphics[width=0.30\textwidth]{imagenes/COPYlogo-utn.eps}\par
    \vspace{1cm}

    \textsc{\Large \@catedra}\par
    \vspace{1.5cm}

    \hrule
    \vspace{0.5cm}
    {\huge\bfseries \@title}\par
    \vspace{0.4cm}
    \textsc{\large \@subtitledoc}
    \vspace{0.5cm}
    \hrule \vspace{\baselineskip}


  \end{center}

  \vspace{2cm}

  {\noindent
    \begin{minipage}[t]{.2\textwidth}
      \raggedright
      \textbf{DOCENTES} \par
      \vspace{\baselineskip}
      %\textbf{CÁTEDRA} \par
      \medskip
      \medskip
      \textbf{COMISIÓN} \par
      \medskip
      \medskip
      \textbf{ALUMNOS} \par
    \end{minipage}%
    \begin{minipage}[t]{.05\textwidth}
      \raggedright
      %\textbf{:} \par
      %~\\
      %~\\
      %\textbf{:} \par
      %\textbf{:} \par
      %\medskip
      %\medskip
      %\textbf{:} \par
    \end{minipage}%
    \begin{minipage}[t]{.55\textwidth}
      \raggedright
      \@professor. \par
      \medskip
      %\@catedra. \par
      \medskip
      \@curso \par
      \medskip
      \medskip
      \@author \par
    \end{minipage}%
    \begin{minipage}[t]{.15\textwidth}
      \raggedright
      \vspace{\baselineskip}
      \vspace{\baselineskip}
      \medskip
      \bigskip
      % \medskip
      \medskip
      \medskip
      \@legajo \par
    \end{minipage}
  }
  \vfill
  \begin{center}
    {\Large\textbf{Córdoba, \@date}} \vspace{\baselineskip}
    %\medskip
    %\textbf{Córdoba, Argentina} \par
    \newpage
  \end{center}
}

\makeatother

\usepackage[nottoc,numbib]{tocbibind}
\usepackage{amsmath}
\usepackage{siunitx}
\usepackage{amssymb}
\usepackage{xfrac}
\usepackage{multirow}
\usepackage{multicol}
\usepackage{float}
\usepackage{circuitikz}
\usepackage{tikz}
\usepackage{pgfplots}
\pgfplotsset{compat=1.18}\instituto{Universidad Tecnológica Nacional\\Facultad Regional Córdoba}
\usetikzlibrary{calc, positioning}



\instituto{Universidad Tecnológica Nacional\\Facultad Regional Córdoba}
\carrera{Ingeniería Electrónica}
\title{Titulo}
\subtitledoc{Subtitulo}
\catedra{Catedra}
\professor{XXXXXXXXXX XXXXXXXX. \par XXXXXXXXXX XXXXXXXX.}
\curso{XRX}
\author{XXXXX XXXXX, XXXXX XXXXX. \par XXXXX XXXXX, XXXXX XXXXX.}
\legajo{XXXXX\par XXXXX}
\footerauthor{XXXXX, XXXXX}
\footerlegajo{XXXXX, XXXXX}
\footercatedra{XXXX}
\date{\today}

\newcommand\plotScale{0.7}

\begin{document} %Maximo 10 paginas
  \maketitle
  \tableofcontents
  \newpage
  \section{Introducción}
  \section{Marco teorico}
  \section{Primera Parte}
  \subsection{Circuito}
  \begin{circuitikz}[american]
    \draw (0,0) to[V, l=$V_{CC}$, invert] ++(0,4)
          to[short] ++(3,0) 
          to[R] ++(0,-2)
          to[biD*] ++(0,-2)
          to[short] (0,0) node[ground]{}
          ;
  \end{circuitikz}
  \subsection{Procedimiento}
  \begin{enumerate}
    \item Armar el circuito seleccionando un correcto valor de R en función del
      datasheet del DIAC.
    \item Variar la tensión de alimentación (V1) desde 0V a 50V según la tabla que se
      observa aquí abajo.
    \item Medir la corriente y caída de tensión en el DIAC.
    \item Invertir los terminales del DIAC y repetir las variaciones y mediciones
      expresadas en el punto 1 y 2
  \end{enumerate}
  \subsection{Calculo de R} %iP = 0.3A (50-14)/0.3  (ST datasheet) DB3
  \subsection{Simulación}
  \begin{figure}[H]
    \begin{center}
      %\includegraphics[width=0.5\textwidth, clip, trim=0.5cm 10cm 0.5cm 8cm]{inc/sim/circuits/CircuitoSegundaParte.pdf}
    \end{center}
    \caption{Circuito simulado}
    \label{fig:CircuitoPrimeraParte}
  \end{figure}
  \begin{center}
    \begin{tikzpicture}[
      ]
      \begin{axis}[
        width=\plotScale\linewidth,
        xlabel=$V_{A_1A_2}$, 
        ylabel=$I$, 
        grid=both,
        thick
        ]
        \addplot[color=red,
        skip coords between index={0}{71},
        ] table [x index=1, y index=3]{./inc/sim/PrimeraParte.txt};
        \addplot[color=blue,
        skip coords between index={69}{142},
        ] table [x index=1, y index=3]{./inc/sim/PrimeraParte.txt};
      \end{axis}
    \end{tikzpicture}
  \end{center}
  \subsection{Experimental}
  \begin{table}[H]
    \begin{center}
      \begin{tabular}{c|c|c}
        $V_{CC}$ &$V_{A_1A_2}$ &$I$  \\
        \hline
        5     &5    &0     \\
        10    &10   &1   uA   \\
        15    &15   &1.5 uA \\
        20    &20   &2   uA   \\
        22    &22   &2.2 uA \\
        24    &24   &2.4 uA \\
        26    &26   &2.6 uA \\
        28    &28   &2.7 uA \\
        30    &30   &3   uA   \\
        32    &32   &3.2 uA \\
        33.4  &27.2 &1.90mA\\
        34    &26.8 &2.24mA\\
        36    &26.1 &3.08mA\\
        38    &25.6 &3.80mA\\
        40    &25.1 &4.10  \\
        42    &24.8 &5.34mA\\
        45    &24.3 &6.41mA \\
        48.7  &23.9 &7.62mA  
      \end{tabular}
    \end{center} 
  \end{table}

  \begin{table}[H] %R = 1000ohm
    \begin{center}
      \begin{tabular}{c|c|c}
        $V_{CC}$ &$V_{A_1A_2}$ &$I$  \\
        \hline
        5     &5    &0     \\
        10    &10   &1    uA\\
        15    &15   &1.5  uA\\
        20    &20   &2    uA\\
        22    &22   &2.2  uA\\
        24    &24   &2.4  uA\\
        26    &26   &2.6  uA\\
        28    &28   &2.7  uA\\
        30    &30   &3    uA\\
        32    &32   &3.2  uA\\
        33.4  &23.2 &10.5 mA\\
        34    &23.2 &11.1 mA\\
        35    &23.0 &12.3 mA\\
        35.5  &22.9 &12.8 mA\\
        36    &22.8 &13.5 mA\\
        37    &22.7 &14.7 mA\\
        38    &22.6 &15.8 mA\\
        40    &22.5 &18.1 mA\\
        42    &22.2 &20.7 mA\\
        45    &22   &23.7 mA \\
        50    &22.1 &29.1 mA \\
        55    &22.3 &34.6 mA \\
        60    &22.5 &40.2 mA \\
        63    &22.6 &43.5 mA \\
      \end{tabular}
    \end{center} 
  \end{table}

 \begin{table}[H] %R = 1000ohm
    \begin{center}
      \begin{tabular}{c|c|c}
        $V_{CC}$ &$V_{A_1A_2}$ &$I$  \\
        \hline
        -5    &-5    &-0     \\
       -10    &-10   &-1    uA\\
       -15    &-15   &-1.5  uA\\
       -20    &-20   &-2    uA\\
       -22    &-22   &-2.2  uA\\
       -24    &-24   &-2.4  uA\\
       -26    &-26   &-2.6  uA\\
       -28    &-28   &-2.7  uA\\
       -30    &-30   &-3    uA\\
       -32    &-32   &-3.2  uA\\
       -33    &-23   &-10.36 mA\\
       -34    &-22.7 &-11.53 mA\\
       -36    &-22.4 &-13.72 mA\\
       -38    &-22.2 &-16.07 mA\\
       -40    &-22.0 &-18.16 mA\\
       -42    &-21.9 &-20.60 mA\\
       -45    &-21.8 &-24.10 mA \\
       -50    &-21.7 &-29.50 mA \\
       -55    &-21.9 &-34.40 mA \\
       -60    &-22.1 &-39.70  mA \\
       -63    &-22.3 &-43.70  mA \\
      \end{tabular}
    \end{center} 
  \end{table}


  \begin{tikzpicture}
    \begin{axis}[
      xlabel = $V_{A_1A_2}$,
      ylabel = $I$,
      grid=major,
      ]
      \addplot[color=red,mark=*] coordinates {
          (5   ,   0   )
          (10  ,   1e-6 )
          (15  ,   1.5e-6)
          (20  ,   2e-6  )
          (22  ,   2.2e-6)
          (24  ,   2.4e-6)
          (26  ,   2.6e-6)
          (28  ,   2.7e-6)
          (30  ,   3e-6  )
          (32  ,   3.2e-6)
          (27.2,   1.90e-3)
          (26.8,   2.24e-3)
          (26.1,   3.08e-3)
          (25.6,   3.80e-3)
          (25.1,   4.10e-3)
          (24.8,   5.34e-3)
          (24.3,   6.41e-3)
          (23.9,   7.62e-3)
        };  
    \end{axis} 
  \end{tikzpicture}


  \begin{tikzpicture}
    \begin{axis}[
      width=0.8\linewidth,
      xlabel = $V_{A_1A_2}$,
      ylabel = $I$,
      grid=major,
      ]
      \addplot[color=red,mark=*, mark size=1.2pt] coordinates {
          (-22.3 ,  -43.70e-3)
          (-22.1 ,  -39.70e-3)
          (-21.9 ,  -34.40e-3)
          (-21.7 ,  -29.50e-3)
          (-21.8 ,  -24.10e-3)
          (-21.9 ,  -20.60e-3)
          (-22.0 ,  -18.16e-3)
          (-22.2 ,  -16.07e-3)
          (-22.4 ,  -13.72e-3)
          (-22.7 ,  -11.53e-3)
          (-23   ,  -10.36e-3)
          (-32   ,  -3.2e-6  )
          (-30   ,  -3  e-6  )
          (-28   ,  -2.7e-6  )
          (-26   ,  -2.6e-6  )
          (-24   ,  -2.4e-6  )
          (-22   ,  -2.2e-6  )
          (-20   ,  -2e-6    )
          (-15   ,  -1.5e-6  )
          (-10   ,  -1e-6    )
          (-5    ,   0       )
          (5     ,   0       )
          (10    ,   1e-6    )
          (15    ,   1.5e-6  )
          (20    ,   2e-6    )
          (22    ,   2.2e-6  )
          (24    ,   2.4e-6  )
          (26    ,   2.6e-6  )
          (28    ,   2.7e-6  )
          (30    ,   3e-6    )
          (32    ,   3.2e-6  )
          (23.2  ,   10.5e-3 )
          (23.2  ,   11.1e-3 )
          (23.0  ,   12.3e-3 )
          (22.9  ,   12.8e-3 )
          (22.8  ,   13.5e-3 )
          (22.7  ,   14.7e-3 )
          (22.6  ,   15.8e-3 )
          (22.5  ,   18.1e-3 )
          (22.2  ,   20.7e-3 )
          (22    ,   23.7e-3 )
          (22.1  ,   29.1e-3 )
          (22.3  ,   34.6e-3 )
          (22.5  ,   40.2e-3 )
          (22.6  ,   43.5e-3 )
        };  
    \end{axis} 
  \end{tikzpicture}
  \section{Segunda Parte}
  \subsection{Circuito}
  \begin{center}
    \begin{circuitikz}[american]
      \draw (0,0) node[vcc](vcc){20V}
            to[R=$R_1$, a=3.3\unit{\kilo\ohm}] ++(0,-1.5)
            to[pR, n=pot, a=5\unit{\kilo\ohm}] ++(0,-4)
            to[short] ++(0,-0.5) node[ground](gnd1){};
      \draw (6,0) node[vcc](VCC){$V_{CC}$}
            to[R=$R_4$, a=4.7\unit{\kilo\ohm}] ++(0,-2.1)
            to[Ty*,n=tr,mirror, v^=$V_{AK}$] (VCC |- pot.wiper)
            to[short,f^=$I_{AK}$] (VCC|-gnd1) node[ground]{};
      \draw (tr.gate) to[nos, mirror, invert] ++(-2,0) 
            to[R=$R_3$, a=4.7\unit{\kilo\ohm}] (pot.wiper);
    \end{circuitikz}
  \end{center}
  \subsection{Procedimiento}
  \begin{enumerate}
    \item Armar el circuito.
    \item Colocar la VCC = 0V.
    \item Cerrar el interruptor Sw.
    \item Variar el potenciómetro de forma de relevar la tabla.
    \item Graficar los valores obtenidos y comparar la curva con la de otro
      componente ya estudiado.
    \item Abrir el interruptor Sw.
    \item Colocar un voltímetro en paralelo a la resistencia de carga (RL) y otro en paralelo al Anodo-Catodo del SCR.
    \item Variar la VCC desde 0V a 600V en pasos de 10V controlando permanentemente lo que sucede en los voltímetros.
    \item Finalizado el ensayo, ¿noto un cambio de comportamiento en el circuito?, ¿En qué valor de tensión?.
    \item Desconectando las alimentaciones de tensión, ¿puede analizar el valor ohmico de la 
      resistencia de carga (RL)?, ¿Qué sucedió?.
    \item Colocar la VG = 0V y cerrar el interruptor Sw.
    \item Colocar VCC = 100V.
    \item Subir lentamente el valor de VG hasta observar un cambio importante en la IAK (Disparo del SCR). 
      Tomar nota del valor de VG e IG que produjo ese disparo del SCR.
    \item Manteniendo el potenciómetro en la posición donde generó el disparo abrir el 
      interruptor Sw y analizar que sucede con la IAK.
    \item Manteniendo el interruptor Sw abierto bajar el valor de VCC en pasos de 10V 
      anotando el valor de IAK para cada caso. Los últimos 10V antes de llegar a cero deben disminuirse en pasos de 1V.
    \item Volver a subir paulatinamente la VCC hasta colocarla nuevamente en 100V
      analizando el comportamiento que tiene la IAK
  \end{enumerate}
  \subsection{Simulación}
  \begin{figure}[H]
    \begin{center}
      %\includegraphics[width=0.5\textwidth, clip, trim=0.5cm 10cm 0.5cm 8cm]{inc/sim/circuits/CircuitoSegundaParte.pdf}
    \end{center}
    \caption{Circuito simulado}
    \label{fig:CircuitoSegundaParte}
  \end{figure}
  \begin{center}
    \begin{tikzpicture}[
      ]
      \begin{axis}[
        width=\plotScale\linewidth,
        xlabel=$V_{D}$, 
        ylabel=$I$, 
        grid=both,
        thick
        ]
        \addplot[color=red,
        %skip coords between index={0}{0},
        ] table [x index=1, y index=2]{./inc/sim/SegundaParte.txt};
      \end{axis}
    \end{tikzpicture}
  \end{center}
  \subsection{Experimental}
\begin{minipage}{0.3\linewidth}
\centering
  \begin{table}[H]
    \begin{center}
      \begin{tabular}{c|c}
        $V_{G}$ &$I_{G}$   \\
        \hline
        0.2 &0   \\
        0.3 &0   \\
        0.4 &0   \\
        0.5 &0   \\
        0.6 &0   \\
        0.7 &0   \\
        0.8 &0   \\
        0.9 &0   \\
        1.0 &0   \\
        1.1 &0   \\
        1.2 &0   \\
        1.3 &0   \\
      \end{tabular}
    \end{center} 
  \end{table}
\end{minipage}
\begin{minipage}{0.7\linewidth}
\centering
\begin{tikzpicture}
    \begin{axis}[
      xlabel = $V_{G}$,
      ylabel = $I_G$,
      grid=major,
      ]
      \addplot[color=red,mark=*] coordinates {
          (5 , 0)
          (10, 0)
          (15, 0)
          (20, 0)
          (22, 0)
          (24, 0)
          (26, 0)
          (28, 0)
          (30  0)
          (32, 0)
          (34, 0)
          (36, 0)
          (38, 0)
          (40, 0)
          (45, 0)
          (50, 0)
        };  
    \end{axis} 
  \end{tikzpicture}
\end{minipage}
\section{Tercera Parte}
\subsection{Circuito}
Para este terecera parte, se trabajo con el mismo circuito que en la parte anterior.
\subsection{Procedimiento}
\begin{enumerate}
  \item Armar el circuito.
  \item Completar los valores de la tabla 1 fijando el valor de IG y variando el valor
        de VCC hasta observar el disparo del SCR. Relevar el valor de VCC, IAK y VAK
        en el cual se produce el disparo para cada caso 
  \item Completar los valores de la tabla 2 fijando el valor de IG y variando el valor
        de VCC. Relevar IAK y VAK para cada caso. Es posible que deba alterar
        algunos valores de resistencia dependiendo del SCR utilizado.
\end{enumerate}
\subsection{Simulación}
  \begin{figure}[H]
    \begin{center}
      %\includegraphics[width=0.5\textwidth, clip, trim=0.5cm 10cm 0.5cm 8cm]{inc/sim/circuits/CircuitoSegundaParte.pdf}
    \end{center}
    \caption{Circuito simulado}
    \label{fig:CircuitoPrimeraParte}
  \end{figure}
  \begin{center}
    \begin{tikzpicture}[
      ]
      \begin{axis}[
        width=\plotScale\linewidth,
        xlabel=$V_{A_1A_2}$, 
        ylabel=$I$, 
        grid=both,
        thick
        ]
        \addplot[color=blue,
        %skip coords between index={69}{142},
        ] table [x index=1, y index=3]{./inc/sim/TerceraParte.txt};
      \end{axis}
    \end{tikzpicture}
  \end{center}
  \subsection{Experimental}
  \begin{table}[H]
    \begin{center}
      \begin{tabular}{c|c|c|c}
        $I_{G}[\unit{\micro\ampere}]$ &$V_{CC}$ &$I_{AK}$ &$V_{AK}$ \\
        \hline
        10 &   &   & \\
        20 &   &   & \\
        20 &   &   & \\
        40 &   &   & \\
        50 &   &   & \\
      \end{tabular}
    \end{center} 
  \end{table}

  \begin{table}[H]
    \begin{center}
      \begin{tabular}{|l|c|c|c|c|c|c|c|c|c|c|}
        \hline
        $I_{G}\rightarrow$ &\multicolumn{2}{c|}{50} &\multicolumn{2}{c|}{40} &\multicolumn{2}{c|}{30} &\multicolumn{2}{c|}{20} &\multicolumn{2}{c|}{10} \\ 
        \hline
        \hline
        $V_{CC}$ &$V_{AK}$ &$I_{AK}$ &$V_{AK}$ &$I_{AK}$ &$V_{AK}$ &$I_{AK}$ &$V_{AK}$ &$I_{AK}$ &$V_{AK}$ &$I_{AK}$\\ 
        \hline
        0    &   &  &  &  &  &  &  &  &  &  \\
        1    &   &  &  &  &  &  &  &  &  &  \\
        2    &   &  &  &  &  &  &  &  &  &  \\
        3    &   &  &  &  &  &  &  &  &  &  \\
        4    &   &  &  &  &  &  &  &  &  &  \\
        5    &   &  &  &  &  &  &  &  &  &  \\
        6    &   &  &  &  &  &  &  &  &  &  \\
        7    &   &  &  &  &  &  &  &  &  &  \\
        8    &   &  &  &  &  &  &  &  &  &  \\
        9    &   &  &  &  &  &  &  &  &  &  \\
        10   &   &  &  &  &  &  &  &  &  &  \\
        11   &   &  &  &  &  &  &  &  &  &  \\
        12   &   &  &  &  &  &  &  &  &  &  \\
        13   &   &  &  &  &  &  &  &  &  &  \\
        14   &   &  &  &  &  &  &  &  &  &  \\
        15   &   &  &  &  &  &  &  &  &  &  \\
        \hline
      \end{tabular}
    \end{center}
  \end{table}



\section{Cuarta Parte}
\subsection{Circuito}
\begin{center}
    \begin{circuitikz}[american]
      \draw (0,0) node[vcc](vcc){20V}
            to[R=$R_1$, a=3.3\unit{\kilo\ohm}] ++(0,-1.5)
            to[pR, n=pot, a=5\unit{\kilo\ohm}] ++(0,-4)
            to[short] ++(0,-0.5) node[ground](gnd1){};
      \draw (6,0) node[vcc](VCC){$V_{CA}$}
            to[R=$R_4$, a=4.7\unit{\kilo\ohm}] ++(0,-2.1)
            to[Ty*,n=tr,mirror, v^=$V_{AK}$] (VCC |- pot.wiper)
            to[R=$R5$,a=100\unit{\ohm},f^=$I_{AK}$] (VCC|-gnd1) node[ground]{};
      \draw (tr.gate) to[nos, mirror, invert] ++(-2,0) 
            to[R=$R_3$, a=4.7\unit{\kilo\ohm}] (pot.wiper);
    \end{circuitikz}
  \end{center}
\subsection{Procedimiento}
\begin{enumerate}
  \item Armar el circuito.
  \item Variar el potenciómetro de forma que VG quede a potencial de cero volts
        cuando conectemos la fuente de alimentación.
  \item Colocar la tensión de alimentación alterna en 50 VCC
  \item Aumentar lentamente la VG observando permanentemente la IG e IA.
        Determinar el momento donde el dispositivo se dispara.
  \item Bajar el valor de VG a cero y observar lo que sucede con la IA.
  \item En función de lo estudiado en el teórico apagar el SCR.
  \item Subir ahora el valor de VCC a 100V y repetir los pasos 4, 5 y 6.
  \item Subir ahora el valor de VCC a 150V y repetir los pasos 4, 5 y 6.
  \item Completar la siguiente tabla de mediciones.
\end{enumerate}
\subsection{Simulación}
  \begin{figure}[H]
    \begin{center}
      %\includegraphics[width=0.5\textwidth, clip, trim=0.5cm 10cm 0.5cm 8cm]{inc/sim/circuits/CircuitoSegundaParte.pdf}
    \end{center}
    \caption{Circuito simulado}
    \label{fig:CircuitoPrimeraParte}
  \end{figure}
  \begin{center}
    \begin{tikzpicture}[
      ]
      \begin{axis}[
        width=\plotScale\linewidth,
        xlabel=$V_{A_1A_2}$, 
        ylabel=$I$, 
        grid=both,
        thick
        ]
        \addplot[color=blue,
        %skip coords between index={69}{142},
        ] table [x index=1, y index=3]{./inc/sim/CuartaParte.txt};
      \end{axis}
    \end{tikzpicture}
  \end{center}
  \subsection{Experimental}

  \end{document}
